% -*- coding: utf-8 -*-
% !TEX program = xelatex
\documentclass[plain,noanswer,medmath]{../../template/jnuexam}
\usepackage{../../template/kysx}

\begin{document}


\examtitle{1999}{3}


\loadproblems{1}{1999P1}%载入试卷一


\makepart{填空题}{本题共5小题，每小题3分，满分15分}


\begin{problem}
设$f(x)$有一个原函数$\frac{\sin x}{x}$，则$\int_{\frac{\pi}{2}}^{\pi}xf'(x)\dx =$ \fillin{}．
\end{problem}


\begin{problem}
$\sum_{n = 1}^{\infty}n\left(\frac{1}{2} \right)^{n - 1} =$ \fillin{}．
\end{problem}


\begin{problem}
设$A = \begin{pmatrix}
  1 & 0 & 1 \\
  0 & 2 & 0 \\
  1 & 0 & 1
\end{pmatrix}$，而$n \ge 2$为整数，则$A^n - 2A^{n - 1} =$ \fillin{}．
\end{problem}


\begin{problem}
在天平上重复称量一重为$a$的物品，假设各次称量结果相互独立且同服从正态分布$N(a, 0.2^2)$．
若以$\overline{X}_n$表示$n$次称量结果的算术平均值，则为使
$$P\left\{\left| \overline{X}_n - a \right| < 0.1 \right\} \ge 0.95,$$
$n$的最小值应不小于自然数 \fillin{}．
\end{problem}


\begin{problem}
设随机变量$X_{ij}\left(i,j = 1,2, \cdots ,n;n \ge 2 \right)$独立同分布，$EX_{ij} = 2$，
则行列式$Y = \begin{vmatrix}
  X_{11} & X_{12} & \cdots & X_{1n} \\
  X_{21} & X_{22} & \cdots & X_{2n} \\
  \vdots & \vdots &        & \vdots \\
  X_{n1} & X_{n2} & \cdots & X_{nn}
\end{vmatrix}$的数学期望$EY =$ \fillin{}．
\end{problem}


\makepart{选择题}{本题共5小题，每小题3分，满分15分}


\useproblem{1}{2}{1}%【同试卷一第二(1)题】


\begin{problem}
设$f(x,y)$连续，且$f(x,y) = xy + \iint_D f(u,v)\du\dv$，
其中$D$是由$y = 0$, $y = x^2$, $x = 1$所围成的区域，则$f(x,y)$等于\pickout{}
\begin{abcd}
\item $xy$．
\item $2xy$．
\item $xy + \frac{1}{8}$．
\item $xy + 1$．
\end{abcd}
\end{problem}


\begin{problem}
设向量$\beta$可由向量组$\alpha_1,\alpha_2, \cdots ,\alpha_m$线性表示，
但不能由向量组(I)$\alpha_1,\alpha_2, \cdots ,\alpha_{m - 1}$线性表示，
记向量组(II)$\alpha_1,\alpha_2, \cdots ,\alpha_{m - 1},\beta$，则 \pickout{}
\begin{abcd}
\item $\alpha_m$不能由(I)线性表示，也不能由(II)线性表示．
\item $\alpha_m$不能由(I)线性表示，但可由(II)线性表示．
\item $\alpha_m$可由(I)线性表示，也可由(II)线性表示．
\item $\alpha_m$可由(I)线性表示，但不可由(II)线性表示．
\end{abcd}
\end{problem}


\begin{problem}
设$A,B$为$n$阶矩阵，且$A$与$B$相似，$E$为$n$阶单位矩阵，则\pickout{}
\begin{abcd}
\item $\lambda E - A = \lambda E - B$．
\item $A$与$B$有相同的特征值和特征向量．
\item $A$与$B$都相似于一个对角矩阵．
\item 对任意常数$t$，$tE - A$与$tE - B$相似．
\end{abcd}
\end{problem}


\begin{problem}
设随机变量$X_i \sim \begin{pmatrix}
           -1 &           0 & 1 \\
  \frac{1}{4} & \frac{1}{2} & \frac{1}{4}
\end{pmatrix}(i = 1,2)$，且满足$P\left\{X_1X_2 = 0 \right\} = 1$，
则$P\left\{X_1 = X_2 \right\}$等于\pickout{}
\begin{abcd}
\item $0$．
\item $\frac{1}{4}$．
\item $\frac{1}{2}$．
\item $1$．
\end{abcd}
\end{problem}


\makepart{}{本题满分6分}

\begin{problem}
曲线$y = \frac{1}{\sqrt{x}}$的切线与$x$轴和$y$轴围成一个图形，记切点的横坐标为$a$，
试求切线方程和这个图形的面积，当切点沿曲线趋于无穷远时，该面积的变换趋势如何？
\end{problem}


\makepart{}{本题满分7分}

\begin{problem}
计算二重积分$\iint_D y\dx\dy$，其中$D$是由直线$x = - 2, y = 0, y = 2$
以及曲线$x = - \sqrt{2y - y^2}$所围成的平面区域．
\end{problem}


\makepart{}{本题满分6分}

\begin{problem}
设生产某种产品必须投入两种要素，$x_1$和$x_2$分别为两要素的投入量，$Q$为产出量；
若生产函数为$Q = 2x_1^{\alpha}x_2^{\beta}$，其中$\alpha , \beta$为正常数，且$\alpha + \beta = 1$．
假设两种要素的价格分别为$p_1$和$p_2$，试问：当产出量为$12$时，两要素各投入多少可以使得投入总费用最小？
\end{problem}


\makepart{}{本题满分6分}

\begin{problem}
设有微分方程$y' - 2y = \varphi(x)$，其中$\varphi (x) = \begin{cases}
  2, & x < 1; \\
  0, & x > 1.
\end{cases}$试求在$\left(-\infty , +\infty \right)$内的连续函数$y = y(x)$，
使之在$\left(-\infty ,1 \right)$和$\left(1, +\infty \right)$内都满足所给方程，且满足条件$y(0) = 0$．
\end{problem}


\makepart{}{本题满分6分}

\begin{problem}
设函数$f(x)$连续，且$\int_0^x tf(2x - t)\dt = \frac{1}{2}\arctan x^2$．
已知$f(1) = 1$，求$\int_1^2 f(x)\dx$的值．
\end{problem}


\makepart{}{本题满分7分}

\begin{problem}
设函数$f(x)$在区间$\left[0,1 \right]$上连续，在$\left(0,1 \right)$内可导，
且$f(0) = f(1) = 0, f\left(\frac{1}{2} \right) = 1$．试证：
\begin{enumerate}
\item 存在$\eta \in \left(\frac{1}{2},1 \right)$，使$f\left(\eta \right) = \eta$；
\item 对任意实数$\lambda$，必存在$\xi \in \left(0,\eta \right)$，
      使得$f'\left(\xi \right) - \lambda \left[f\left(\xi \right) - \xi \right] = 1$．
\end{enumerate}
\end{problem}


\makepart{}{本题满分9分}%【同试卷一第十题】

\useproblem{1}{10}{0}


\makepart{}{本题满分7分}

\begin{problem}
设$A$为$m \times n$实矩阵，$E$为$n$阶单位矩阵．已知矩阵$B = \lambda E + A^T A$，
试证：当$\lambda > 0$时，矩阵$B$为正定矩阵．
\end{problem}


\makepart{}{本题满分9分}

\begin{problem}
假设二维随机变量$\left(X, Y \right)$在矩形
$G = \left\{\left. \left(x, y \right) \right|0 \le x \le 2, 0 \le y \le 1 \right\}$上服从均匀分布．记
$$U = \begin{cases}
  0, & X \le Y; \\
  1, & X > Y.
\end{cases}\quad V = \begin{cases}
  0, & X \le 2Y; \\
  1, & X > 2Y.
\end{cases}$$
\begin{enumerate*}
\item 求$U$和$V$的联合分布；
\item 求$U$和$V$的相关系数$r$．
\end{enumerate*}
\end{problem}


\makepart{}{本题满分7分}

\begin{problem}
设$X_1,X_2, \ldots ,X_9$是来自正态总体$X$的简单随机样本，
$Y_1 = \frac{1}{6}\left(X_1 + X_2 + \cdots + X_6 \right)$，%%
$Y_2 = \frac{1}{3}\left(X_7 + X_8 + X_9 \right)$，%%
$S^2 = \frac{1}{2}\sum_{i = 7}^9 \left(X_i - Y_2 \right)^2$，%%
$Z = \frac{\sqrt{2} \left(Y_1 - Y_2 \right)}{S}$，证明统计量$Z$服从自由度为$2$的$t$分布．
\end{problem}


\end{document}
